\documentclass[11pt]{article}
\usepackage{amsmath,amssymb,amsthm}
\usepackage[margin=1.1in]{geometry}
\usepackage{hyperref}
\newtheorem{theorem}{Theorem}
\newtheorem{lemma}[theorem]{Lemma}
\newtheorem{remark}[theorem]{Remark}
\newtheorem{problem}[theorem]{Problem}
\newcommand{\C}{\mathbb{C}}

\title{Low-degree classification of torus-equivariant Keller maps\\ in dimension three
\thanks{Draft. Computational classification, unconditional as of 2026-07-21
(Level-3 certification, Section~5, discharged; see the status box after the
title for artifacts). All certificates, solver transcripts, and reproduction
files accompany this note.}}
\author{Anonymous\thanks{Pseudonymous deposit. Correspondence: \texttt{minustsquared@gmail.com}. The authors' identity is committed to by a SHA-256 hash published in the accompanying \texttt{COMMITMENT.txt}.}}
\date{July 2026 --- draft v0.1}

\begin{document}
\maketitle

\begin{abstract}
Motivated by the recent counterexample to the Jacobian conjecture in dimension
three announced by Alp\"oge and credited to the AI system Fable
\cite{alpoge2026}, which is homogeneous for the $\C^{*}$-action with weights
$(1,-1,-2)$ and has polynomial degree seven, we classify all Keller maps
compatible with the same torus action in polynomial degree at most four:
every $(1,-1,-2)$-equivariant Keller map of degree at most four is a
polynomial automorphism (Theorem~\ref{thm:main}). The classification was
first established computationally in conditional form and became
unconditional on 2026-07-21 once the decomposition was certified exhaustive
in two independent computer algebra systems (Section~\ref{sec:pending}). The
question of whether the same holds in degrees five and six --- equivalently,
whether degree seven is minimal within this symmetry class --- is posed as an
open computational problem, and the degree-five system (13 unknowns, 17
equations) is supplied.
\end{abstract}

\section{Introduction}

A polynomial map $F\colon\C^{n}\to\C^{n}$ is a \emph{Keller map} if
$\det \operatorname{Jac} F$ is a nonzero constant. The Jacobian conjecture
\cite{keller1939,bcw1982} asserted that every Keller map is a polynomial
automorphism; it is now false for $n\ge 3$ by the explicit counterexample of
\cite{alpoge2026}. The verification paper \cite{ulam2026} checks, in exact
arithmetic, that the counterexample's Jacobian determinant is a nonzero
constant and that the map is not injective (an explicit multi-point fiber is
exhibited), and derives a binary-cubic description of the fibers, of generic
cardinality three; the present note relies only on these verified facts
about the example. The
counterexample has polynomial degree seven and is homogeneous with respect to
the grading $\operatorname{wt}(x,y,z)=(1,-1,-2)$, with component weights
$(-2,-1,1)$.

A natural first question about the mechanism is whether it is available in
lower polynomial degree within the same symmetry class. This note answers that
question negatively for degree at most four (Theorem~\ref{thm:main}), first in
conditional form and, as of 2026-07-21, unconditionally once the reduction to
a single certification task (Section~\ref{sec:pending}) was discharged.

We emphasize the distinction between three evidentiary levels used throughout:
\begin{itemize}
\item[\textbf{L1}] \emph{Proved by hand:} the equivariant ansatz and the
forced-normalization lemma (Lemma~\ref{lem:normal}).
\item[\textbf{L2}] \emph{Exactly machine-certified:} Jacobian-constancy
identities, explicit inverse formulas, composition identities, denominator
conditions, and the $b_0=0$ stratum, with certificates as symbolic objects over
exact fields (no floating point).
\item[\textbf{L3}] \emph{Discharged 2026-07-21:} certified exhaustive
decomposition of the Keller ideal, cross-checked in two independent systems
(Section~\ref{sec:pending}). The main theorem is unconditional.
\end{itemize}

\section{The equivariant family and forced normalization}

Fix the $\C^{*}$-action $t\cdot(x,y,z)=(tx,t^{-1}y,t^{-2}z)$ and require
$F=(F_1,F_2,F_3)$ to have components of weights $(-2,-1,1)$, the weights of the
known counterexample. For a degree bound $D$, write each $F_i$ as a linear
combination of the monomials of the prescribed weight and total degree at most
$D$.

\begin{lemma}[Forced normalization]\label{lem:normal}
The only degree-one monomials of weights $-2$, $-1$, $1$ are $z$, $y$, $x$
respectively, so the linear part of any such $F$ is anti-diagonal,
$(\alpha z,\beta y,\gamma x)$, and $\det\operatorname{Jac}F(0)=-\alpha\beta\gamma$.
If $F$ is Keller then $\alpha\beta\gamma\neq0$, and $F$ may be normalized to
$\alpha=\beta=\gamma=1$ by equivariant diagonal linear changes of source and
target. Hence the normalized family considered below contains every
$(1,-1,-2)$-equivariant Keller map of degree at most $D$ up to such changes.
\end{lemma}

\begin{proof}
Weight bookkeeping gives the monomial claim; the constant $\det\operatorname{Jac}F$
equals its value at the origin, which is the determinant of the linear part.
Diagonal linear maps $\operatorname{diag}(s,t,u)$ commute with the torus action,
so pre- and post-composition by such maps preserves the equivariance class;
the six available scalings (with a one-dimensional redundancy from the torus
itself) normalize the three anchor coefficients to $1$.
\end{proof}

For $D=4$ the normalized family is explicitly
\[
F=\bigl(z+a_0y^2+a_1xyz+a_2xy^3+a_3x^2z^2,\;
y+b_0xz+b_1xy^2+b_2x^2yz,\;
x+c_0x^2y+c_1x^3z\bigr),
\]
nine unknown coefficients; imposing
constancy of $\det\operatorname{Jac}F$ yields $11$ polynomial equations over
$\mathbb{Q}$ (the \emph{Keller ideal}). Note that for weight-homogeneous $F$
the determinant is automatically of weight zero, i.e.\ a polynomial in the
invariants $xy$ and $x^{2}z$; this is the structural reason the constancy
condition is a manageable system.

\section{Certification methodology (L2)}\label{sec:method}

The Keller ideal was decomposed with SymPy~1.14.0 (Python~3.12.3) via
Gr\"obner-based triangular decomposition (\texttt{sympy.solve}). The raw
case-split returns $27$ branches, several re-derived along redundant
elimination paths; deduplicating by solution locus leaves $5$ distinct
parametrized families. (Correction, 2026-07-21: earlier circulated drafts of
this note quoted the raw, undeduplicated branch count; no family or
certificate was lost by deduplicating, only redundant re-derivations of the
same locus.) The completeness of this decomposition was item L3, discharged
2026-07-21 (Section~\ref{sec:pending}). For each family, with free parameters
kept \emph{symbolic}:
\begin{enumerate}
\item the substituted determinant is expanded and verified to be a constant
identically in $(x,y,z)$ (it equals $-1$ on every family);
\item a lexicographic Gr\"obner basis of
$\langle F_1-p,\,F_2-q,\,F_3-r\rangle$ ($x>y>z$) is computed over the exact
function field $\mathbb{Q}(\text{parameters})(p,q,r)$, the unique solution
$(x,y,z)=(X,Y,Z)(p,q,r)$ is extracted, and the composition identities
$F_i(X,Y,Z)=p,q,r$ are verified by exact cancellation;
\item all irreducible factors of all denominators of $X,Y,Z$ are computed. If
every factor is a nonzero constant, the inverse is valid for \emph{every}
parameter specialization in the family's domain (\emph{pointwise} certificate);
otherwise the factors are recorded as exceptional.
\end{enumerate}
A rational inverse with polynomial re-composition to the identity makes the map
birational; a birational Keller map is a polynomial automorphism. Indeed, a
Keller map is \'etale, hence quasi-finite, and a separated, quasi-finite,
birational morphism onto a normal variety is an open immersion by Zariski's
main theorem \cite{mumford}; applied to $F\colon\C^{n}\to\C^{n}$ this makes
$F$ injective, and an injective polynomial self-map of $\C^{n}$ is an
automorphism with polynomial inverse \cite{rudin}. When the certificate is
pointwise, this conclusion holds at every parameter value, not only generically.

Outcome: $4$ of $5$ distinct families certify pointwise. The remaining $1$
certifies with the single exceptional factor $b_0$ (the coefficient of $xz$
in $F_2$), hence pointwise on the constructible piece $b_0\neq0$, which is the
entire stated domain of that parametrization. The complementary constructible
piece is handled by decomposing the ideal augmented by $b_0$: after the same
deduplication it yields $1$ distinct family, which certifies pointwise.
(Second correction, 2026-07-21: this stratum's decomposition initially
appeared to yield $2$ distinct families, for $7$ total; one of those two
turned out to be an exact duplicate of a \emph{main}-branch family --- the
main decomposition, not told to assume $b_0=0$, had independently landed on
a branch already sitting there. Deduplicating across both files, not only
within each, leaves $6$ distinct families total, $5$ from the main piece and
$1$ genuinely new from the stratum.) All certificates are supplied as
machine-readable JSON with the full inverse formulas. An independent
verifier --- a separate program that reads only the static certificates and
checks the determinant identities, both compositions
$F\circ H=\operatorname{id}$ and $H\circ F=\operatorname{id}$, and the
recorded denominator factors, performing no solving --- passes on all $6$
distinct families. The ideal-theoretic certification
(Section~\ref{sec:pending}) separates the two loci as the saturation
$I:b_0^{\infty}$ and the sum $I+\langle b_0\rangle$ and requires their
components to be covered by the certified families -- carried out and
discharged 2026-07-21.

\section{Main result}

\begin{theorem}[Computational classification]\label{thm:main}
Every $(1,-1,-2)$-equivariant Keller map $\C^3\to\C^3$ of polynomial degree at
most four is a polynomial automorphism.
\end{theorem}

The Keller ideal $I$ has exactly three minimal primes (Singular~4.4.1
\texttt{minAssGTZ}, independently reproduced in Macaulay2~1.26.06
\texttt{decompose}). Minimal primes suffice because the theorem is a
statement about \emph{points}: the set of normalized equivariant Keller maps
is the affine variety $V(I)$, and set-theoretically
$V(I)=V(P_1)\cup V(P_2)\cup V(P_3)$ over the minimal primes, irrespective of
any embedded associated primes, which contribute no additional points (no
radicality claim about $I$ is made or needed). The relation between the six
certified families and the three components is as follows
(\texttt{out/certify.log}): families $0$, $2$, and $5$ are dense charts of
the three components --- component $1\leftrightarrow$ family $5$,
$2\leftrightarrow2$, $3\leftrightarrow0$ --- with family $5$ the
$b_0$-stratum family, so one entire component lies inside $\{b_0=0\}$; the
remaining families $1$, $3$, $4$ are constructible subfamilies contained in
these closures, which arose as separate branches of the triangular
decomposition and are retained with their own certificates. Every point of
every component carries an explicit polynomial inverse with both
compositions verified: the pointwise certificates are valid on their entire
loci, and the single non-pointwise chart (family $0$, valid on $b_0\neq0$)
is complemented by the $b_0=0$ stratum decomposition
(Section~\ref{sec:method}).

\begin{remark}
The theorem should be read together with Lemma~\ref{lem:normal}: the
normalization is not a hypothesis but a reduction, so the conclusion covers
the full equivariant class up to equivariant diagonal linear changes.
\end{remark}

\begin{remark}[Certification history]
Earlier drafts stated Theorem~\ref{thm:main} conditionally on the
exhaustiveness of the computed decomposition; that certification is now
discharged (Section~\ref{sec:pending}). The certification process surfaced
and corrected two family-count bookkeeping errors and two computer-algebra
usage bugs, none affecting the mathematical content. The full audit trail
--- the family-count correction history and the specific Singular and
Macaulay2 pitfalls, recorded so that others do not lose time to them --- is
kept in the accompanying repository (\texttt{HANDOFF.md}) rather than here.
\end{remark}

\section{Certification history (L3, discharged 2026-07-21)}\label{sec:pending}

The single previously-uncertified step was exhaustiveness of the two
decompositions. The accompanying files \texttt{certify\_exhaustiveness.sing}
(Singular, \texttt{minAssGTZ} plus containment checks against the 6 certified
family ideals) and \texttt{deg4\_decomposition.m2} (Macaulay2,
\texttt{decompose}) compute the minimal associated primes of the Keller ideal
over $\mathbb{Q}$. Both agree on 3 minimal primes; every one is covered by a
certified family, and every certified family lies inside the Keller variety
(both containment directions checked, see \texttt{out/certify.log}). The
theorem is therefore unconditional: no omitted component remains that could
contain maps outside the certified families.

\section{Open problems}

\begin{problem}
Decide degrees five and six. The degree-five Keller system ($13$ unknowns,
$17$ equations, supplied in \texttt{deg5\_system.txt}) exceeded the pure-Python
Gr\"obner engine used here; an exact decomposition in Singular is running as
of this writing (2026-07-21) but has not finished, apparently limited by
intermediate coefficient growth rather than the size of the search space --
three independent modular runs (primes 32003, 32009, 65003) each find exactly
$2$ minimal primes with identical structure, and three exact families
reconstructed from them (one splitting into two further branches) each
independently verify: $\det=-1$ and an explicit polynomial inverse with
verified composition, by the same method as the certified degree-4 families
(\texttt{04-scripts/reconstruct\_deg5\_candidates.py}). \textbf{This is not a
classification result for degree five} -- the modular component count is a
heuristic preview, not a characteristic-zero certificate, and exhaustiveness
of these families has not been established by any method. A negative answer
through degree six would establish degree seven as minimal within this
symmetry class and invite a conceptual explanation, possibly via the
binary-cubic fiber mechanism of \cite{ulam2026}.
\end{problem}

\begin{problem}
Enumerate all weight systems $(w_1,w_2,w_3)$ admitting nontrivial equivariant
Keller families in low degree, and run the same classification across them.
\end{problem}

\begin{problem}
Replace the branch-by-branch computation with a structural proof, e.g.\ a
triangular factorization theorem for low-degree equivariant Keller maps.
\end{problem}

\section*{Acknowledgment}

The formulation and the exploratory symbolic computations were carried out
with the assistance of Claude Fable (Anthropic). All results are certified
by machine-checkable artifacts in exact arithmetic --- certificates, solver
transcripts, and an independent solving-free verifier in the accompanying
repository --- independent of any AI system's reasoning and reproducible
from the supplied materials. Correspondence: \texttt{minustsquared@gmail.com}.

\begin{thebibliography}{9}
\bibitem{alpoge2026} L.~Alp\"oge, X post, July 20, 2026,
\url{https://x.com/__alpoge__/status/2079028340955197566}.
\bibitem{ulam2026} A counterexample to the Jacobian conjecture, verification
paper, July 20, 2026, \url{https://www.ulam.ai/research/jacobian.pdf}.
\bibitem{keller1939} O.-H.~Keller, Ganze Cremona-Transformationen,
Monatsh.\ Math.\ Phys.\ 47 (1939), 299--306.
\bibitem{bcw1982} H.~Bass, E.~H.~Connell, D.~Wright, The Jacobian conjecture:
reduction of degree and formal expansion of the inverse, Bull.\ Amer.\ Math.\
Soc.\ 7 (1982), 287--330.
\bibitem{rudin} W.~Rudin, Injective polynomial maps are automorphisms,
Amer.\ Math.\ Monthly 102 (1995), 540--543.
\bibitem{mumford} D.~Mumford, \emph{The Red Book of Varieties and Schemes},
Lecture Notes in Mathematics 1358, Springer, 1988 (Zariski's main theorem).
\end{thebibliography}

\end{document}
